% !TeX root = ../main.tex
\section*{Обозначения и соглашения}
\addcontentsline{toc}{section}{Обозначения и соглашения}
Набраны и проверены решения из файла
\texttt{дз1\_чат\_кураторов.pdf}. Нумерация сохранена по рукописи:
«С.\,$k$:$n$» означает задачу $n$ из параграфа $k$ сборника под редакцией
В.\,К.\,Романко; «Ф.» обозначает номера задачника
А.\,Ф.\,Филиппова; «Д.» --- дополнительные задачи без указания сборника.
Если источник в рукописи явно не указан, это оговорено в тексте.
Программа курса из другого приложенного файла не заменяет состав этой домашней работы.

Все решения вещественные. Буквы $C,C_1,C_2$ обозначают произвольные
постоянные, если ограничения не указаны отдельно. После переобозначения
постоянной её прежнее численное значение не сохраняется.
Производная $y'$ берётся по $x$; при смене независимой переменной это
указывается явно. Кубический корень понимается как вещественный,
в частности $x^{4/3}=(\sqrt[3]{x})^4$.

Решение $y=y(x)$ рассматривается на интервале, где оно имеет необходимые
производные и удовлетворяет исходному уравнению. Для уравнений
$P\dd x+Q\dd y=0$ отдельно допускаются параметрические интегральные
кривые, в том числе вертикальные. Прямая $x=0$ не является графиком
функции $y(x)$: эти два смысла решения не следует смешивать.

При делении на выражение, которое может обращаться в нуль, исключённые
случаи проверяются в исходном уравнении. Неявная формула $F(x,y)=C$
задаёт график лишь на подходящих гладких ветвях. Условие $F_y\ne0$
обеспечивает возможность локально выразить $y$ через $x$.
Уравнения первого порядка допускают $C^1$-склейки, если в точке
соединения совпадают значения функций и первые производные.

Утверждения о существовании и единственности в разделе III являются
\emph{локальными}: гладкость правой части во всей плоскости сама по себе
не исключает ухода решения на бесконечность за конечное время.
Основные исправления и соответствие страницам исходника приведены в приложении.
